\section{Predicate logic translation}\label{sec:predicate_logic_translation}

\paragraph{First-order logic to higher-order logic}

\begin{algorithm}[First-order signature to higher-order signature]\label{alg:fol_signature_to_hol_signature}
  We will describe a procedure for converting a \hyperref[def:fol_signature]{first-order signature} \( \Sigma \) to a \hyperref[def:hol_signature]{higher-order signature} \( \Sigma' \).

  \begin{thmenum}
    \thmitem{alg:fol_signature_to_hol_signature/sort} The only sort in \( \Sigma' \) is the type of individuals, \( \syn\iota \).
    \thmitem{alg:fol_signature_to_hol_signature/function} For every \( n \)-ary function symbol \( f \) in \( \Sigma \), we introduce into \( \Sigma' \) a nonlogical constant with the same name and with type
    \begin{equation*}
      \underbrace{\syn\iota \synimplies \cdots \synimplies \syn\iota}_{n \T*{times}} \synimplies \syn\iota.
    \end{equation*}

    \thmitem{alg:fol_signature_to_hol_signature/predicate} For every \( n \)-ary predicate symbol \( p \) in \( \Sigma \), we introduce into \( \Sigma' \) a nonlogical constant with the same name and with type
    \begin{equation*}
      \underbrace{\syn\iota \synimplies \cdots \synimplies \syn\iota}_{n \T*{times}} \synimplies \syn\omicron.
    \end{equation*}
  \end{thmenum}
\end{algorithm}
\begin{comments}
  \item This algorithm can be found as \identifier{math.lambda_.hol.fol_signature_to_hol_signature} in \cite{notebook:code}.
\end{comments}

\begin{algorithm}[First-order term to higher-order expression]\label{alg:fol_term_to_hol_expression}
  Fix a \hyperref[def:fol_term]{first-order term} \( T \) over the \hyperref[def:fol_signature]{first-order signature} \( \Sigma \). Fix also a \hyperref[def:fol_structure]{first-order structure} \( \mscrX = (X, I) \) and a \hyperref[def:fol_variable_assignment]{variable assignment} \( v \) into \( \mscrX \).

  Via recursion on \( T \), we will translate \( T \) into a \hyperref[def:hol_expression]{well-formed expression} \( T' \):
  \begin{thmenum}
    \thmitem{alg:fol_term_to_hol_expression/variable} If \( T \) is a variable, say \( x \), simply let \( T' = x \).
    \thmitem{alg:fol_term_to_hol_expression/application} If \( T = f(S_1, \ldots, S_n) \), let \( S_k' \) be the translation of \( S_k \) and let
    \begin{equation*}
      T \coloneqq (\ldots ((f S_1) S_2) \ldots S_n).
    \end{equation*}
  \end{thmenum}
\end{algorithm}
\begin{comments}
  \item This translation is justified in \cref{thm:fol_to_hol_translation}.
  \item This algorithm can be found as \identifier{math.lambda_.hol.fol_term_to_hol_expression} in \cite{notebook:code}.
\end{comments}

\begin{algorithm}[First-order formula to higher-order expression]\label{alg:fol_formula_to_hol_expression}
  Fix a \hyperref[def:fol_formula]{first-order formula} \( \varphi \) over the \hyperref[def:fol_signature]{first-order signature} \( \Sigma \).

  Via recursion on \( \varphi \), we will translate \( \varphi \) into a \hyperref[def:hol_expression]{well-formed expression} \( \varphi' \):
  \begin{thmenum}
    \thmitem{alg:fol_formula_to_hol_expression/verum} If \( \varphi = \syntop \), simply let \( \varphi' \coloneqq \synH_\top \).
    \thmitem{alg:fol_formula_to_hol_expression/falsum} Similarly, if \( \varphi = \synbot \), simply let \( \varphi' \coloneqq \synH_\bot \).
    \thmitem{alg:fol_formula_to_hol_expression/application} If \( \varphi = p(\psi_1, \ldots, \psi_n) \), let \( \psi'_k \) be the translation of \( \psi_k \) and let
    \begin{equation*}
      \varphi' \coloneqq (\ldots ((p \psi'_1) \psi'_2) \ldots \psi'_n).
    \end{equation*}

    \thmitem{alg:fol_formula_to_hol_expression/connective} If \( \varphi = \psi \syncirc \theta \) , let \( \psi' \) and \( \theta' \) be the translations of \( \psi \) and \( \theta \), and let
    \begin{equation*}
      \varphi' \coloneqq (\synH_{\syncirc} \psi') \theta'.
    \end{equation*}

    \thmitem{alg:fol_formula_to_hol_expression/quantifier} If \( \varphi = \quantifier Q x \psi \) , let \( \psi' \) be the translations of \( \psi \), and let
    \begin{equation*}
      \varphi' \coloneqq \synH_Q (\qabs {x^{\syn\iota}} \psi').
    \end{equation*}
  \end{thmenum}
\end{algorithm}
\begin{comments}
  \item This translation is justified in \cref{thm:fol_to_hol_translation}.
  \item This algorithm can be found as \identifier{math.lambda_.hol.fol_formula_to_hol_expression} in \cite{notebook:code}.
\end{comments}

\begin{algorithm}[First-order structure to higher-order structure]\label{alg:fol_structure_to_hol_structure}
  Fix a \hyperref[def:fol_structure]{first-order structure} \( \mscrX = (X, I) \) over the \hyperref[def:fol_signature]{first-order signature} \( \Sigma \).

  we will translate \( \mscrX \) into a \hi{standard} \hyperref[def:hol_structure]{higher-order structure} \( \mscrX' = (X', I') \) as follows:
  \begin{thmenum}
    \thmitem{alg:fol_signature_to_hol_signature/frame} We define the frame \( X' = \seq{ X'_\tau }_{\tau \in \op*{\BbbQ_{\Sigma'}}} \) as follows:
    \begin{thmenum}
      \thmitem{alg:fol_signature_to_hol_signature/frame/individuals} For the type of individuals, we let \( X'_{\syn\iota} \) be the first-order universe \( X \).

      \thmitem{alg:fol_signature_to_hol_signature/frame/arrow} For the arrow type \( \tau \synimplies \sigma \), we define the universe \( X'_{\tau \synimplies \sigma} \) as the set of all functions from \( X'_\tau \) to \( X'_\sigma \).
    \end{thmenum}

    \thmitem{alg:fol_signature_to_hol_signature/interpretation} We define interpretation \( I' \) as follows: for every \( n \)-ary function or predicate symbol \( s \) in \( \Sigma \), let \( I'(s) \) be \( I_s \), where
    \begin{equation*}
      I_{s,a_1,\ldots,a_k} \coloneqq \begin{cases}
        I(s)(a_1, \ldots, a_n),                       &k = n, \\
        a_{k+1} \mapsto I_{s,a_1,\ldots,a_k,a_{k+1}}, &k < n.
      \end{cases}
    \end{equation*}
  \end{thmenum}
\end{algorithm}
\begin{comments}
  \item It seems reasonable to define \( X'_{\tau \synimplies \sigma} \) as the set of only those functions from \( X'_\tau \) to \( X'_\sigma \) expressible in first-order logic. This would unfortunately leave \( X'_{\tau \synimplies \sigma} \) empty in many cases, for example if \( \tau = \syn\omicron \).

  \item This algorithm can be found as \identifier{math.lambda_.hol.fol_structure_to_hol_structure} in \cite{notebook:code}.
\end{comments}

\begin{proposition}\label{thm:fol_to_hol_translation}
  Fix a \hyperref[def:fol_structure]{first-order structure} \( \mscrX = (X, I) \) over the \hyperref[def:fol_signature]{first-order signature} \( \Sigma \). Fix also a \hyperref[def:fol_variable_assignment]{first-order variable assignment} \( v \) into \( \mscrX \).

  Let \( \Sigma' \) be the signature translated using \fullref{alg:fol_signature_to_hol_signature}, let \( \mscrX' \) be the structure translated using \fullref{alg:fol_structure_to_hol_structure} and let \( v' \) be any \hyperref[def:hol_variable_assignment]{higher-order assignment} that coincides with \( v \) on individuals.

  \begin{thmenum}
    \thmitem{thm:fol_to_hol_translation/term} For any \hyperref[def:fol_term]{first-order term} \( T \), for the translation \( T' \), obtained via \fullref{alg:fol_term_to_hol_expression}, we have
    \begin{equation*}
      \Bracks{T'}_{\mscrX'}^{v'} = \Bracks{T}_\mscrX^v.
    \end{equation*}

    \thmitem{thm:fol_to_hol_translation/formula} Similarly, for any \hyperref[def:fol_formula]{first-order formula} \( \varphi \), for the translation \( \varphi' \), obtained via \fullref{alg:fol_formula_to_hol_expression}, we have
    \begin{equation*}
      \Bracks{\varphi'}_{\mscrX'}^{v'} = \Bracks{\varphi}_\mscrX^v.
    \end{equation*}
  \end{thmenum}
\end{proposition}
\begin{proof}
  \SubProofOf{thm:fol_to_hol_translation/term} We proceed by \fullref{thm:induction_on_abstract_syntax} on \( T \) to show that \( \Bracks{T'}_{\mscrX'}^{v'} = \Bracks{T}_\mscrX^v \).
  \begin{itemize}
    \item If \( T \) is a variable, say \( x \), then \( v'(x^{\syn\iota}) = v(x) \) by assumption.
    \item If \( T = f(S_1, \ldots, S_n) \), where the inductive hypothesis holds for \( S_1, \ldots, S_n \), then
    \begin{balign*}
      \Bracks{T'}_{\mscrX'}^{v'}
      &=
      \Bracks{(\ldots ((f S'_1) S'_2) \ldots S'_n)}_{\mscrX'}^{v'}
      \reloset {\eqref{eq:alg:hol_denotation/app}} = \\ &=
      \Bracks{f}_{\mscrX'}^{v'} \parens[\big]{ \Bracks{S'_1}_{\mscrX'}^{v'} } \ldots \parens[\big]{ \Bracks{S'_n}_{\mscrX'}^{v'} }
      \reloset {\T{ind.}} = \\ &=
      \Bracks{f}_{\mscrX'}^{v'} \parens[\big]{ \Bracks{S_1}_\mscrX^v } \ldots \parens[\big]{ \Bracks{S_n}_\mscrX^v }
      \reloset {\eqref{eq:alg:hol_denotation/abs}} = \\ &=
      I'(f) \parens[\big]{ \Bracks{S_1}_\mscrX^v } \ldots \parens[\big]{ \Bracks{S_n}_\mscrX^v }
      \reloset {\ref{alg:fol_signature_to_hol_signature/interpretation}} = \\ &=
      I(f) \parens[\big]{ \Bracks{S_1}_\mscrX^v, \ldots, \Bracks{S_n}_\mscrX^v }
      = \\ &=
      \Bracks{T}_\mscrX^v.
    \end{balign*}
  \end{itemize}

  \SubProofOf{thm:fol_to_hol_translation/formula} We can use induction on the formula \( \varphi \), similarly to \cref{thm:fol_to_hol_translation/term}, to show that \( \Bracks{\varphi'}_{\mscrX'}^{v'} = \Bracks{\varphi}_\mscrX^v \).
\end{proof}

\paragraph{Higher-order logic to first-order logic}

\begin{algorithm}\label{alg:higher_order_to_first_order_logic}
\end{algorithm}

\begin{algorithm}\label{alg:many_sorted_to_one_sorted_theory}
\end{algorithm}
